\begin{abstract}
This paper studies adult income prediction as a binary classification problem on the public Adult dataset. The task is to predict whether an individual's annual income exceeds \$50K from 14 demographic and employment-related attributes. We compare two substantially different machine learning methods covered in the module: Random Forest and Multi-Layer Perceptron (MLP). Both models are trained and evaluated under the same preprocessing pipeline, stratified 70/15/15 data split, and validation-based model selection protocol. On the held-out test set, Random Forest achieved higher Accuracy (0.8622), Precision (0.7812), and ROC-AUC (0.9141), while MLP achieved higher Recall (0.6269) and F1-score (0.6765). In practical terms, the MLP recovered more high-income cases, whereas Random Forest made fewer false positive predictions. MLP was also the more efficient method, training faster and requiring far less memory and storage. The main finding is therefore not that one model dominates the other, but that the preferred method depends on deployment priorities: Random Forest is the safer choice when conservative positive predictions and simpler interpretation matter most, while MLP is more attractive when balanced retrieval and low computational cost are more important.
\end{abstract}
